% !TeX program = xelatex
% ============================================================
%  新讲义起手式 · tlnote v1.0（2026-11-09 起启用，用户选定 A 方案）
%  新建讲义请复制本文件；不要复制 11/04–11/09 那 30 份旧 .tex 的 preamble。
%
%  编译：xelatex 两遍（目录 / 书签 / 交叉引用才完整）
%        xelatex -interaction=nonstopmode -halt-on-error <本文件>
%  选项：[color] 克制彩色（默认）｜[gray] 灰度打印
%  安装：D:\Run_Tao\Homework\install.ps1；kpsewhich tlnote.sty 可验证
%
%  ★ 命令契约：tlnote.sty 已提供 geometry / titlesec / tcolorbox /
%    fancyhdr / hyperref / tikz / pgfplots / enumitem / booktabs /
%    longtable / tabularx / caption / pifont 及全部页眉页脚与色彩。
%    本文件【不要】再 \usepackage{geometry}、\usepackage{hyperref}、
%    \pagestyle{fancy}、\renewcommand{\headrulewidth}{}、
%    \newtcolorbox{keypoint} 等——重复加载会选项冲突，重复定义会报错。
%
%  ★ 禁止重复定义（会报 already defined，或静默改写正文）：
%    \abs \ee \D \dd \blank \ans \R \eps \N \Z \ue \dif \ds \floor \norm
%    \key \tagline \kw \yes \no \src \ver \card \rowh \q \Q \qi
%    其中 \blank 用【可选参数】\blank[1.2cm]，不是 \blank{1.2cm}。
%
%  ★ 存量 30 份讲义（11/04–11/09）保持 ctexart + 手写 preamble，不迁移：
%    合订本 370 页 / 阅读包 61、74 页 / 书签 / 绝对页码索引绑定原排版。
% ============================================================
\documentclass[11pt,a4paper]{ctexart}
\usepackage[color]{tlnote}

\title{新讲义标题}
\author{陶理}
\noteheader{高等数学 I · 期中复习}
\docdate{2026-11-09}          % 固定日期；禁用 \date{\today}

% 取 \@title / \@author / \@date 的别名（用户文档里 @ 不是字母）
\makeatletter
\newcommand{\showtitle}{\@title}
\newcommand{\showauthor}{\@author}
\newcommand{\showdate}{\@date}
\makeatother

\begin{document}

\begin{center}
  {\LARGE\bfseries\color{tlblue}\showtitle\par}
  \vspace{0.35em}
  {\normalsize\showauthor\quad\showdate\par}
\end{center}
\vspace{-0.3em}
\noindent{\color{tlrule}\rule{\textwidth}{0.5pt}}

% 长讲义再取消下一行的注释；短讲义保持注释可收在 1 页。
% 书签由 \section 经 hyperref 生成，不依赖 \tableofcontents。
% {\small\tableofcontents}

\section{方法卡}

\begin{keypoint}[要点盒]
层级靠字号与间距，颜色只作辅助；盒子只承载明确语义。
数学简写直接用包内命令：$\R$、$\Z$、$\abs{x}$、$\D x$、$\dd{y}{x}$、$\eps$、$\ee^{x}$、$\norm{f}$。
\end{keypoint}

\section{题目与解答}

\qi{例 1}{2024 期中 A 卷 Q3}{求 $\lim\limits_{x\to0}\dfrac{\sin x}{x}$。}
\begin{solution}
由重要极限得 $\ans{1}$。填空用可选参数：\blank[1.2cm]。
\end{solution}

\begin{warnbox}[易错点]
填空命令 \texttt{blank} 用方括号 \texttt{[1.2cm]}；写成花括号会把长度当正文印出来。
在正文里提到命令名时写 \texttt{xxx}，不要裸写带参数的宏。
\key{先定语义，再选命令，不要另起一套宏。}
\end{warnbox}

\begin{methodbox}[一题多解]
\begin{subproblem} 夹逼定理。 \end{subproblem}
\begin{subproblem} 洛必达法则。 \end{subproblem}
\end{methodbox}

\section{性态表（列表写法）}

\begin{center}
\begin{tabular}{c|c|c|c}
\toprule
区间 & $y'$ & $y''$ & 性态 \\
\midrule
\rowh $(-\infty,0)$ & $+$ & $-$ & 增、凸 \\
\rowh $(0,+\infty)$ & $-$ & $+$ & 减、凹 \\
\bottomrule
\end{tabular}
\end{center}

\begin{notebox}[注]
\yes\ 结论经复核；\no\ 未验证的推断不写进讲义。\src{题型专题 §0}
\end{notebox}

\ver{xelatex 两遍；缺字 0；本文件仅为版式起手式，不含练习题，不计入题池。}
\tagline{新建讲义：复制本文件，改标题与 docdate 日期字段，正文直接写，不要重写 preamble。}

\end{document}
